\documentclass[reprint]{revtex4-2}
\usepackage{amsmath,amssymb,graphicx,hyperref,booktabs}
\begin{document}

\title{Information Capacity and the Structure of Physical Law: \\
Why Quantum, Classical, and Gravity are Phases of Finite Information}

\author{Huang Zhongchang}
\affiliation{Independent Researcher, Nanjing, China}

\begin{abstract}
We show that three definitional truths about information---distinguishable states exist (A1), capacity is bounded (A2), and interactions are local (A3)---are sufficient to explain the structural form of the quantum-classical separation, the arrow of time, and the Newtonian gravitational field equation. The central object is the $q$-field, which tracks the vacant information capacity fraction. We prove that the information flux $J = -\kappa\nabla(\ln q)$ is the unique form satisfying Shannon's uniqueness theorem for additive information measures and the physical constraint that $q=1$ (vacuum) is stable. The continuum $q$-field obeys $\partial_t q = \kappa\nabla^2(\ln q)$, whose static limit is $\nabla^2(\ln q)=0$, structurally identical to the Laplace equation for the Newtonian gravitational potential. The $q\to 0$ regime constitutes a self-organized critical state; the $q\approx 1/2$ regime is the entropic attractor where quantum coherence persists. All qualitative conclusions---dual-domain asymptotic states, Laplace-type static limit, $1/r$ gravitational decay---belong to a universality class $\Phi$ and are independent of the specific flux form. The framework makes a falsifiable prediction: $\ln q(r)$ is harmonic in vacuum, yielding $q(r)=e^{-GM/rc^2}$ in spherical symmetry. This implies black holes lack a sharp event horizon ($q>0$ for all $r>0$), in contrast to general relativity. The classical world is not the quantum world's graveyard. It is the quantum world's engine---a self-organized critical state whose static geometry is what we measure as gravity.
\end{abstract}

\section{Introduction}

\subsection{The question}

Why does the universe appear quantum at microscopic scales and classical at macroscopic scales? Standard decoherence theory~\cite{Zurek2003,Schlosshauer2007} answers \emph{how} this occurs: environmental monitoring selects pointer states, and off-diagonal density matrix elements decay. This answer is correct but incomplete. It does not address why the separation happens in this direction---classicality grows, quantum coherence recedes, never the reverse---or why the gravitational field equation shares the mathematical structure of an information diffusion equation in its static limit.

This paper proposes that the quantum-classical separation and the structure of gravity are not independent facts. They are different manifestations of a single principle: \emph{information is finite, and its flow under capacity constraints generates the structures we observe as physical law.}

\subsection{What this paper does}

We establish a chain of reasoning from definitional truths about information to structural features of physical law:

\begin{enumerate}
\item \textbf{Definitions.} To exist is to be distinguishable (A1). To interact is to transfer information within finite bounds (A2). Interactions are local (A3). These are not axioms to be proved---they are operational definitions of what ``existence'' and ``interaction'' mean.

\item \textbf{Flux theorem.} From A1--A3 and Shannon's uniqueness theorem for additive information measures, the information flux between cells takes the form $J = -\kappa\nabla(\ln q)$, where $q$ is the vacant capacity fraction. This is the unique form satisfying the physical constraint that $q=1$ (vacuum) is stable.

\item \textbf{Continuum $q$-field.} The continuum limit yields $\partial_t q = \kappa\nabla^2(\ln q)$. The $q\to 0$ regime is a fast-diffusion singularity---a self-organized critical state. The $q\approx 1/2$ regime is the entropic attractor.

\item \textbf{Universality.} All qualitative conclusions belong to a universality class $\Phi$: any flux function with $\phi(0^+)=+\infty$ produces dual-domain asymptotic states, a Laplace-type static limit, and $1/r$ gravitational decay.

\item \textbf{Gravity.} The static limit $\nabla^2(\ln q)=0$ is structurally identical to the Laplace equation for the Newtonian gravitational potential. The spherical solution $q(r)=e^{-GM/rc^2}$ yields the weak-field correspondence $q\approx 1-GM/rc^2$.

\item \textbf{Testable prediction.} DGF predicts black holes lack a sharp event horizon ($q>0$ for $r>0$), in contrast to general relativity. This can be tested with gravitational wave ringdown and Event Horizon Telescope observations.
\end{enumerate}

\section{Framework}

\subsection{Definitions}

\textbf{A1 (Distinguishability).} For any two states $s\neq s'$, their images under dynamics satisfy $f(s)\neq f(s')$. This is the operational content of ``existence'': to be is to be distinguishable from what is not.

\textbf{A2 (Bounded capacity).} The state space $X=\{0,1\}^N$ is finite, with $|X|=2^N$. Each cell stores 1 bit. Infinite capacity would make ``interaction'' ill-defined---there would be no boundary separating system from environment.

\textbf{A3 (Locality).} The dynamics $f:X\to X$ is $c$-local with respect to some cell decomposition: each output cell depends on at most $c$ input cells, with $c\ll N$. Locality is not derived from A1--A2 (Theorem~OA-1 proves it cannot be~\cite{A4_analysis}); it is acknowledged as a physically motivated input encoding finite propagation speed.

\subsection{The $q$ parameter}

Define the vacant capacity fraction:
\begin{equation}
q = \frac{1}{N}\sum_{i=1}^N (1-s_i) \in [0,1].
\end{equation}
$q=1$: all cells vacant (maximal capacity, fully quantum). $q=0$: all cells occupied (zero capacity, classical). The $q$ parameter admits a precise mapping to standard quantum mechanics: $q(\rho) = C_{\ell_1}(\rho)/\|\rho\|_1$, the normalized $\ell_1$-norm of coherence in the pointer basis~\cite{Baumgratz2014}.

\subsection{Structural theorems}

\begin{theorem}[Bijectivity]
An injective $f:X\to X$ on a finite set is bijective.
\end{theorem}
\begin{proof}
Pigeonhole principle: $|X|$ finite, $f$ injective $\Rightarrow$ $f$ surjective.
\end{proof}

This is the structural conservation law: the total number of distinguishable states is preserved. However, no \emph{specific} macroscopic quantity need be conserved:

\begin{theorem}[Non-conservation]
A1--A2 do not imply $\sum_i s_i$ is constant. Counterexample for $N=2$: $f(00)=01, f(01)=10, f(10)=11, f(11)=00$.
\end{theorem}

\begin{theorem}[Entropy identity]
\label{thm:entropy}
Under coarse-graining $\pi:X\to[0,1]$, $S[P']-S[P] = \mathbb{E}_{q'}[D[\rho'_{q'}||\mu_{q'}]] - \mathbb{E}_q[D[\rho_q||\mu_q]]$.
\end{theorem}
This identity decomposes Shannon entropy change into coarse-grained and intra-fiber contributions. When intra-fiber distributions are initially uniform, entropy is bounded below by its initial value. Entropy can decrease between consecutive steps while remaining above $S[P_0]$~\cite{INSPECTOR_round1}.

\begin{theorem}[Kac recurrence]
The expected recurrence time for a subsystem of $N_S$ cells is $\langle T_{\rm rec}\rangle = 2^{N_S}$, independent of environment size $N_E$.
\end{theorem}
This separates relaxation (controlled by $N_E$) from recurrence (controlled by $N_S$). ``Irreversibility'' is fast relaxation relative to observation time.

\begin{theorem}[Reflux bound]
When total excitation is conserved, $P_{\rm reflux} \le T\cdot(q_S/(1-q_S))^\delta\cdot((1-q_E)/q_E)^\delta$.
\end{theorem}
This bound provides a capacity constraint on quantum dynamics complementary to standard spectral density constraints.

\section{The Flux Theorem}

\subsection{Problem statement}

The discrete dynamics $f$ determines how information moves between cells. In the continuum limit, this is encoded in the flux function $J(q_i,q_j)$. What is its form?

\subsection{Constraints}

Any physically admissible flux must satisfy:
\begin{itemize}
\item \textbf{C1 (Directionality):} $J_{i\to j}>0 \iff q_i < q_j$ (information flows from full to empty).
\item \textbf{C2 (Antisymmetry):} $J_{i\to j} = -J_{j\to i}$.
\item \textbf{C3 (Locality):} $J$ depends only on $q_i,q_j$.
\item \textbf{C4 (Capacity constraint):} $\lim_{q\to 0}D(q)=\infty$ (saturated cells cannot absorb).
\item \textbf{C5 (Vacuum stability):} $D(1)<\infty$ (vacuum is stable against perturbations).
\end{itemize}

\subsection{Shannon uniqueness and the self-information potential}

C1--C3 and the irrotationality condition $\oint J=0$ imply the additive separable form $J_{i\to j} = \Phi(q_j)-\Phi(q_i)$ for some potential $\Phi$~\cite{flux_classification}. The potential $\Phi(q)$ must be an information measure: it quantifies the ``pressure'' exerted by a cell at vacancy $q$.

Shannon's uniqueness theorem~\cite{Shannon1948} states that the only additive information measure continuous in its argument is $I(p) = -K\log p$. Setting $p=q$ (the probability that a randomly selected cell is vacant):
\begin{equation}
\Phi(q) = -\ln q,
\end{equation}
up to an additive constant and multiplicative factor.

The flux follows immediately:
\begin{equation}
\boxed{J_{i\to j} = \ln\frac{q_j}{q_i}},
\end{equation}
or in continuum notation, $J = -\kappa\nabla(\ln q)$ where $\kappa$ is a diffusivity constant. The effective diffusivity is $D(q)=\kappa/q$.

\subsection{Why not other forms?}

The Shannon entropy gradient gives $D(q)\propto 1/[q(1-q)]$, which diverges at \emph{both} $q\to 0$ and $q\to 1$. The $q\to 1$ divergence renders the vacuum unstable---any perturbation propagates instantaneously, contradicting the observed stability of the quantum vacuum. The self-information form has $D(1)=\kappa$, finite---vacuum perturbations propagate at finite speed.

Scale invariance arguments~\cite{flux_derivation} give $D(q)\propto 1/q^2$, which diverges more strongly at $q\to 0$. This belongs to the same universality class (Theorem~\ref{thm:universality}) and produces identical qualitative conclusions. The self-information form is preferred by (a)~Shannon uniqueness, (b)~minimal divergence satisfying C4, and (c)~vacuum stability.

\subsection{Universality}

\begin{theorem}[Universality class $\Phi$]
\label{thm:universality}
Let $\Phi$ be the class of flux functions $\phi(q)$ satisfying: $\phi$ is smooth on $(0,1]$, $\phi'(q)<0$, $\phi(0^+)=+\infty$, and $\phi(1)$ is finite. Then for any $\phi\in\Phi$, the qualitative conclusions---dual-domain asymptotic states, Laplace-type static limit, and $1/r$ gravitational decay---are universal, independent of the specific $\phi$.
\end{theorem}
\begin{proof}[Proof sketch]
The static limit $\partial_t q = -\nabla^2\phi(q) = 0$ gives $\nabla^2\phi(q)=0$ for any $\phi$, which is always a second-order elliptic equation. The Green's function in $d=3$ decays as $1/r$ independently of $\phi$. The dual-domain structure follows from $\phi(0^+)=+\infty$ (forcing $q\to 0$ cores) combined with the entropic preference for $q\approx 1/2$. Full proof in Ref.~\cite{flux_universality}.
\end{proof}

\section{Continuum $q$-Field}

\subsection{Equation of motion}

From the self-information flux $J=-\kappa\nabla(\ln q)$ and continuity $\partial_t q = -\nabla\cdot J$:
\begin{equation}
\boxed{\partial_t q = \kappa\,\nabla^2(\ln q)}.
\label{eq:qfield}
\end{equation}

Equivalent forms: $\partial_t q = \kappa\nabla\cdot(q^{-1}\nabla q)$ with $D(q)=\kappa/q$.

\subsection{The $q\to 0$ singularity}

As $q\to 0$, $D(q)\to\infty$. This is a fast-diffusion singularity---a \emph{self-organized critical state}. Cells are nearly saturated; any incoming perturbation finds almost no buffer capacity; cascaded overflow propagates with near-zero inter-step waiting time. Each step respects the maximum propagation speed $c$; the apparent instantaneous response arises from the vanishing interval between steps, analogous to sandpile avalanches at the angle of repose~\cite{Bak1987}.

\subsection{The $q\approx 1/2$ attractor}

The mixing entropy per cell $s(q)=-q\log q-(1-q)\log(1-q)$ attains its unique maximum at $q=1/2$. The state-space volume is maximal there: $\binom{N}{N/2}\sim 2^N/\sqrt{\pi N/2}$. Statistical mechanics drives the background toward $q\approx 1/2$.

Two competing forces shape the $q$-field:
\begin{itemize}
\item \textbf{Diffusion pressure}: the $1/q$ nonlinearity amplifies $q$-gradients, creating $q\approx 0$ cores.
\item \textbf{Entropic pressure}: the statistical preference drives the background toward $q\approx 1/2$.
\end{itemize}

The competition produces a two-domain asymptotic state: compact $q\approx 0$ cores (classical objects) embedded in a $q\approx 1/2$ sea (quantum background).

\subsection{The separation mechanism}

Start from $q\approx 1$ (the initial quantum universe). Any spatial inhomogeneity creates a $q$-gradient. The $1/q$ nonlinearity amplifies this gradient: regions of lower $q$ lose information faster, driving them toward $q\approx 0$. These cores become critically sensitive---unable to absorb information, they immediately re-emit it. Surrounding regions, maintained near $q\approx 1/2$ by entropic pressure, receive this outflow while retaining quantum coherence.

The separation is not two worlds diverging. It is one system developing two phases under a single dynamics.

\section{Gravity from the Critical State}

\subsection{Static limit}

In the static limit $\partial_t q=0$, Eq.~(\ref{eq:qfield}) reduces to:
\begin{equation}
\boxed{\nabla^2(\ln q) = 0}.
\label{eq:laplace}
\end{equation}

This is the Laplace equation. The information potential $\psi\equiv -\ln q$ satisfies $\nabla^2\psi=0$---the same equation satisfied by the Newtonian gravitational potential in vacuum, $\nabla^2\Phi=0$.

With source terms (mass as locked information):
\begin{equation}
\nabla^2\psi = -\kappa\rho_m,
\end{equation}
structurally identical to the Poisson equation $\nabla^2\Phi = 4\pi G\rho$. The \emph{form} of the equation is derived; the \emph{value} of the coupling $\kappa$ is calibrated to $4\pi G/c^2$ by matching to Newtonian gravity.

\subsection{Spherical solution}

Imposing spherical symmetry and the boundary condition $q(\infty)=1$:
\begin{equation}
\boxed{q(r) = e^{-GM/rc^2}},
\label{eq:spherical}
\end{equation}
valid in the weak-field regime. Expanding: $q(r)\approx 1 - GM/rc^2$, yielding the Newtonian correspondence $\ln q = \Phi/c^2$.

\subsection{Scale comparison}

\begin{table}[h]
\caption{$q$ at surface for systems spanning the mass range.}
\label{tab:scales}
\begin{tabular}{lccc}
\toprule
System & $M$ [kg] & $R$ [m] & $q(R)$ \\
\midrule
Earth & $6.0\times 10^{24}$ & $6.4\times 10^6$ & $0.9999999993$ \\
Sun & $2.0\times 10^{30}$ & $7.0\times 10^8$ & $0.9999979$ \\
Neutron star & $1.4M_\odot$ & $10^4$ & $0.83$ \\
Black hole (at $r_s$) & $10M_\odot$ & $3\times 10^4$ & $e^{-1/2}\approx 0.607$ \\
\bottomrule
\end{tabular}
\end{table}

Earth is an extraordinarily quantum object: less than $10^{-7}$\% of cells occupied. Yet this minuscule deviation, integrated over Earth's volume and amplified by $1/q$ nonlinearity, produces $g=9.8$~m/s$^2$.

\subsection{Black holes: no sharp event horizon}

The solution $q(r)=e^{-GM/rc^2}$ satisfies $q(r)>0$ for all $r>0$. There is no finite radius at which $q=0$---no sharp event horizon. The approach to $q=0$ is asymptotic as $r\to 0$.

This contrasts with general relativity, where $g_{00}=0$ at the Schwarzschild radius $r_s=2GM/c^2$. In DGF, ``black holes'' are objects with $q$ exponentially suppressed at small $r$, but information is never strictly trapped. The S1 reflux bound gives $P_{\rm reflux}\propto e^{-N_S}$, non-zero in principle.

\textbf{Observational signatures:} (1)~Gravitational wave ringdown: the absence of a sharp horizon modifies the quasi-normal mode spectrum. (2)~EHT black hole images: the photon orbit and shadow size differ from GR predictions. We emphasize that these signatures require a hyperbolic extension of Eq.~(\ref{eq:qfield}); the current parabolic form describes the static Newtonian limit. The predictions are presented as qualitative targets for future development.

\subsection{What is derived and what is calibrated}

\textbf{Derived structure:} The Laplace/Poisson form of the field equation, the $1/r$ gravitational decay, Gauss's law structure, and the absence of a sharp event horizon all follow from Eq.~(\ref{eq:qfield}).

\textbf{Calibrated constants:} $G$ is fixed by matching to Newtonian gravity. The mass-information conversion factor $\eta$ [kg/bit] is calibrated, not derived.

\textbf{Open:} The current equation is parabolic. A hyperbolic extension is required for gravitational waves. The strong-field regime ($r\lesssim GM/c^2$) requires a discrete treatment.

\section{Physical Picture}

\subsection{The information flow narrative}

A pure quantum universe begins with $q\approx 1$ (all capacity vacant). Initial inhomogeneities create $q$-gradients. Information flows from low-$q$ to high-$q$ regions---statistically, not as a law, driven by the surprisal gradient $\nabla(-\ln q)$. Regions approaching $q\approx 0$ enter self-organized criticality, becoming powerful sources of $q$-field gradients. These critical regions are what we call mass. Their static geometry, extending to infinity as $1/r$, is what we measure as gravity. Surrounding regions at $q\approx 1/2$ receive this outflow while retaining quantum coherence.

The classical world is not the quantum world's graveyard. It is its engine---a self-organized critical state perpetually pumping information outward.

\subsection{Time}

Overflow events establish directed edges between cells. The set of all such edges forms a directed acyclic graph. The partial order on this DAG is time: event $A$ precedes $B$ if $A$ is a causal ancestor of $B$ in the overflow record. Before the first overflow, no edges exist---no partial order, no time. This is the Euclidean phase. The first overflow creates the first causal relation---the origin of time. The arrow of time is the statistical impossibility of reversing the overflow DAG when the capacity ratio is large.

\subsection{Relation to other frameworks}

\textbf{Zurek's quantum Darwinism~\cite{Zurek2009}:} Standard decoherence requires an external environment. DGF shows that capacity asymmetry alone---without external monitoring---suffices to produce preferred bases and classicality. The two are complementary: quantum Darwinism describes the mechanism; DGF describes the condition under which it is inevitable.

\textbf{Bassani \& Magueijo~\cite{Bassani2025}:} Their framework addresses \emph{how} fixed laws emerge through natural selection. DGF addresses \emph{why} the laws take the structural form they do. BM provides the mechanism; DGF provides the architecture. The two are complementary.

\textbf{Penrose/Di\'osi~\cite{Penrose1996,Diosi1989}:} These models link decoherence to gravitational self-energy. DGF links gravity to information capacity exhaustion. Both associate gravity with the quantum-classical transition, but through different mechanisms.

\section{Experimental Contact}

The framework makes one clean prediction that distinguishes it from general relativity: \textbf{black holes lack a sharp event horizon}. The $q$-field $q(r)=e^{-GM/rc^2}$ satisfies $q>0$ everywhere, implying information is never strictly trapped.

\textbf{Gravitational wave ringdown:} The quasi-normal mode spectrum of a DGF ``black hole'' differs from GR because no sharp boundary condition exists at $r_s$. Current LIGO/Virgo ringdown measurements~\cite{LIGO2021} are consistent with GR but have not ruled out soft-boundary alternatives.

\textbf{Event Horizon Telescope:} The photon orbit in a $q$-field background differs from GR at small $r$. The shadow size is predicted to be smaller than the GR value of $\sim 5.2\,GM/c^2$, approaching $\sim 2.7\,GM/c^2$ in a refractive-index model. Current EHT measurements of M87*~\cite{EHT2019} and Sgr A*~\cite{EHT2022} constrain but do not exclude such deviations.

We emphasize that quantitative predictions require: (i)~a hyperbolic extension of the $q$-field equation for propagating gravitational waves; (ii)~an independent measurement protocol for $q$ in astrophysical contexts. These are open problems, not solved results.

The S1 reflux bound (Theorem~5) provides a complementary laboratory-scale test: at zero temperature, DGF permits finite information reflux $P_{\rm reflux}\le q_S/q_E$, while standard Lindblad theory predicts strictly zero. Current circuit QED technology is approaching the required sensitivity~\cite{circuit_qed_prediction}.

\section{Honest Boundaries}

\subsection{What the framework explains}

\begin{enumerate}
\item \textbf{Why quantum and classical behavior coexist.} They are two phases of the same $q$-field at different capacity ratios.
\item \textbf{Why the separation is directional.} The $1/q$ nonlinearity amplifies toward $q\approx 0$; no symmetric mechanism amplifies toward $q\approx 1$.
\item \textbf{Why gravity has the Laplace/Poisson form.} It is the static limit of the $q$-field equation.
\item \textbf{Why time has a direction.} The overflow DAG encodes a partial order; reversing it is statistically impossible at large capacity ratios.
\end{enumerate}

\subsection{What the framework does not explain}

\begin{enumerate}
\item The numerical values of $c$, $\hbar$, $G$.
\item Why $d=3$ spatial dimensions.
\item Quantum field theory, particle physics, the standard model.
\item The specific scale of the quantum-classical boundary (why at $\sim 10^{-10}$--$10^{-6}$~m).
\item Gravitational waves (requires hyperbolic extension).
\item Strong-field gravity (requires discrete treatment).
\end{enumerate}

\subsection{Assumption inventory}

\begin{table}[h]
\caption{Complete assumption inventory.}
\begin{tabular}{lll}
\toprule
Label & Content & Status \\
\midrule
A1 & Distinguishable states exist & Definition \\
A2 & Capacity bounded & Definition \\
A3 & Locality & Physical input \\
C5 & Vacuum stability ($D(1)<\infty$) & Physical input \\
$\kappa$ & Flux normalization & Calibrated to $G$ \\
$\eta$ & Mass-information conversion & Calibrated to $G$ \\
$d$=3 & Spatial dimension & Observation \\
\bottomrule
\end{tabular}
\end{table}

\section{Conclusion}

We have shown that three definitions and one physical input---distinguishability, bounded capacity, locality, and vacuum stability---are sufficient to explain the structural form of the quantum-classical separation, the arrow of time, and the Newtonian gravitational field equation. The framework's central result is the $q$-field equation $\partial_t q = \kappa\nabla^2(\ln q)$, whose static limit is the Laplace equation for the gravitational potential. The flux form $J=-\kappa\nabla(\ln q)$ is selected by Shannon's uniqueness theorem as the only additive information measure satisfying vacuum stability.

The deepest insight is perhaps the simplest: the classical world is not the quantum world's graveyard. It is its engine---a self-organized critical state whose static geometry is what we measure as gravity. Quantum coherence is not fragile. It is the default, sustained by the entropic preference for $q\approx 1/2$. The separation of scales is not a mystery. It is a phase coexistence driven by the competition between diffusive pressure and entropic attraction.

We do not claim to have derived the universe. We claim to have shown that certain structural features of our physical world---features that appear contingent in standard treatments---are necessary consequences of what it means for information to exist, to be finite, and to flow.

\begin{thebibliography}{99}

\bibitem{Zurek2003} W.H.~Zurek, Rev.~Mod.~Phys.~\textbf{75}, 715 (2003).
\bibitem{Schlosshauer2007} M.~Schlosshauer, \emph{Decoherence and the Quantum-to-Classical Transition} (Springer, 2007).
\bibitem{Zurek2009} W.H.~Zurek, Nat.~Phys.~\textbf{5}, 181 (2009).
\bibitem{Baumgratz2014} T.~Baumgratz, M.~Cramer, and M.B.~Plenio, Phys.~Rev.~Lett.~\textbf{113}, 140401 (2014).
\bibitem{Bak1987} P.~Bak, C.~Tang, and K.~Wiesenfeld, Phys.~Rev.~Lett.~\textbf{59}, 381 (1987).
\bibitem{Shannon1948} C.E.~Shannon, Bell Syst.~Tech.~J.~\textbf{27}, 379 (1948).
\bibitem{Penrose1996} R.~Penrose, Gen.~Rel.~Grav.~\textbf{28}, 581 (1996).
\bibitem{Diosi1989} L.~Di\'osi, Phys.~Rev.~A~\textbf{40}, 1165 (1989).
\bibitem{Bassani2025} P.M.~Bassani and J.~Magueijo, Phys.~Rev.~D~\textbf{111}, 103529 (2025).
\bibitem{Cover2006} T.M.~Cover and J.A.~Thomas, \emph{Elements of Information Theory} (Wiley, 2006).
\bibitem{Chiribella2011} G.~Chiribella, G.M.~D'Ariano, and P.~Perinotti, Phys.~Rev.~A~\textbf{84}, 012311 (2011).
\bibitem{Wheeler1990} J.A.~Wheeler, in \emph{Complexity, Entropy, and the Physics of Information} (1990).

\bibitem{flux_derivation} Expert~1, \emph{Flux derivation from first principles}, LP31-InfoCapacity (2026).
\bibitem{flux_universality} Expert~2, \emph{Universality of flux conclusions}, LP31-InfoCapacity (2026).
\bibitem{flux_classification} Expert~3, \emph{Classification of flux forms}, LP31-InfoCapacity (2026).
\bibitem{A4_analysis} A4 analysis, \emph{Status of the locality axiom}, LP31-InfoCapacity (2026).
\bibitem{INSPECTOR_round1} INSPECTOR, \emph{Round 1 verification}, LP31-InfoCapacity (2026).
\bibitem{circuit_qed_prediction} B博士, \emph{Circuit QED prediction for S1}, LP31-InfoCapacity (2026).

\bibitem{LIGO2021} LIGO/Virgo Collaboration, Phys.~Rev.~D~\textbf{103}, 122002 (2021).
\bibitem{EHT2019} EHT Collaboration, Astrophys.~J.~\textbf{875}, L1 (2019).
\bibitem{EHT2022} EHT Collaboration, Astrophys.~J.~\textbf{930}, L12 (2022).

\end{thebibliography}

\end{document}
